\section{Euler packaging, coprimality and the analytic boundary}
\label{rc:details}
For $p\notin T_0$, set $z_p=\psi_u(p)(Np)^{-s}$ and
$t_p=z_p/(1+z_p)$. On $\operatorname{Re}s>1$, deleting an additional
surviving prime divides the squarefree Euler product by $1+z_p$.
Consequently the exact retained expression is
\[
 J_{u,U}(s)=F_{u,1}(s)B_{u,U}(s),\qquad
 B_{u,U}(s)=\sum_{d\ \mathrm{squarefree},\ (d,T_0)=1}
                 \beta_U(d)\prod_{p\mid d}t_p.
\]
Original zeros enforce $(d,u)=1$ even when a sixth power drops from
its primitive conductor. This factorization cannot restore a forbidden
$d$ row by replacing a zero sentinel with an unramified primitive value.
The equality continues wherever both sides are holomorphic.

Without the cutoffs, $\beta_\infty(d)=\mu(d)2^{\omega(d)}$ on
squarefree ideals. Absolute convergence justifies
\[
 B_{u,\infty}=\prod_{p\notin T_0}(1-2t_p)
 =\prod_{p\notin T_0}\frac{1-z_p}{1+z_p},\qquad
 F_{u,1}B_{u,\infty}=\prod_{p\notin T_0}(1-z_p)
                   =\frac1{L^{T_0}(s,\psi_u)}.
\]
Theorem~\ref{rc:recombination} gives a quantitative rate rather than an
unsupported exchange of this limit with a critical-strip contour.

\paragraph{A local counterexample to the naive square.}
In a one-prime projection at a good Eisenstein prime of norm $7$, take
the principal local phase and $s=1$. Then $z=1/7$, $t=1/8$ and,
for $U\ge7$, the exact polynomial is $B_U=1-2t=3/4$.
The unconstrained square $(1-t)^2=49/64$ differs by $1/64$:
it admits the forbidden pair $a=b=p$. After multiplying by
$F=1+z=8/7$ it gives $7/8$ instead of the correct $6/7$.
This is an actual local Euler calculation, not an infinite-series or
mean-square counterexample. It identifies the exact coprimality condition
that a proposed mollifier-square simplification must preserve.

\paragraph{The exact intermediate coefficient band.}
For a squarefree $m$ with $U<Nm\le U^2$, one has $A_U(m)=-\mu(m)$.
Indeed the unrestricted disjoint-assignment sum is $\mu(m)$. If $Na>U$,
every $b\mid m/a$ automatically has $Nb\le Nm/Na<U$, and
$\sum_{b\mid m/a}\mu(b)$ vanishes unless $a=m$. Hence the excluded
$a$ assignments contribute exactly $\mu(m)$; the excluded $b$
assignments contribute another $\mu(m)$. They cannot overlap because
$NaNb>U^2\ge Nm$ contradicts $ab\mid m$. Subtraction proves the claim.
Together with the small-norm band, this gives $A_U=\mu$ for $Nm\le U$
and $A_U=-\mu$ for $U<Nm\le U^2$. The original smoothed columns have
$Nm\ge D\ge U^2$, so the new band does not justify replacing their
coefficients by $-\mu(m)$: generally they lie beyond it.

\begin{theorem}[A stronger cutoff convergence would require zero-freeness]
\label{rc:boundary}
Fix a nonprincipal primitive $\psi_u$. If these actual $J_{u,U}$
converge locally uniformly on $\operatorname{Re}s>1/2$ as
$U\longrightarrow\infty$, then $L(s,\psi_u)$ has no zero there.
\end{theorem}
\begin{proof}
Classical abelian Hecke continuation makes the numerator of
$F_{u,1}=L^{T_0}(s,\psi_u)/L^{T_0}(2s,\psi_{u,2})$ holomorphic.
Its denominator is nonzero on $\operatorname{Re}s>1/2$ by the
absolute Euler product at $2s$. Each $B_{u,U}$ is a finite sum with
nonzero denominators $1+z_p$ on $\operatorname{Re}s>0$.
Thus every $J_{u,U}$ is holomorphic on the asserted half-plane.
A locally uniform limit $G$ is holomorphic there. On
$\operatorname{Re}s>1$, the already proved convergence gives
$L^{T_0}(s,\psi_u)G(s)=1$. The identity theorem extends this equality
to the connected half-plane, precluding a zero of $L^{T_0}$.
The finitely deleted factors $1-\psi_u(p)(Np)^{-s}$ do not vanish
when $\operatorname{Re}s>0$, so the same conclusion holds for the
primitive $L$-function.
\end{proof}
This is an implication, not a proof that the stronger convergence holds,
and it does not assert a converse. In particular, a family-averaged energy
bound is not silently promoted to locally uniform convergence for each
individual character. Principal rows have a pole at $1$ and were excluded
from this proposition. General Artin holomorphy is not an input.

\paragraph{Reproducible exact diagnostics.}
The new source \path{examples/arithmetic_cutoff/exact_recombination.py}
uses actual split/inert prime ideals and raw sextic symbols in
$\mathbb Z[\omega]$, with exact rational arithmetic. It checks the direct
masked sum, coefficient regrouping, Euler packaging, both finite full-cutoff
limits and original zeros, including third/sixth-power rows. Its overlapping-
support negative controls are not tests of an infinite large sieve or of
zero location. The analytic convergence argument above uses ordinary ideal
Euler products and the explicitly stated Hecke continuation theorem; no
new Lean declaration or complete GRH proof is claimed.
